% Fichier complété par tools/complete_missing_corrections.py.
% Source : RDM/RDM-01-Cohesion/526_RdM/526_RdM.tex
% Statut : rédigé (3 question(s)).
\begin{corrigebox}[Corrigé rédigé]
\CorrigeQuestion{1}
On coupe la poutre à l'abscisse $x$ et on isole la partie la plus simple.
            Le PFS donne $\{T_{coh}\}=-\sum\{T_{ext}\}$ au centre de la section. On
            obtient les composantes $N(x)$, $T_y(x)$, $T_z(x)$,
            $M_t(x)$, $M_{fy}(x)$ et $M_{fz}(x)$ tronçon par tronçon.
\CorrigeQuestion{2}
Une composante $N$ correspond à traction/compression, $T_y,T_z$ au
            cisaillement, $M_t$ à la torsion et $M_{fy},M_{fz}$ à la flexion. Les
            composantes non nulles du torseur précédent donnent directement les
            sollicitations combinées.
\CorrigeQuestion{3}
Chaque diagramme est tracé avec les expressions par morceaux. Les sauts
            correspondent aux forces ou couples concentrés; la pente de l'effort tranchant
            vaut l'opposé de la charge répartie et la pente du moment fléchissant est
            l'effort tranchant. Les trois figures \texttt{cor\_01.jpg}, \texttt{cor\_02.jpg} et
            \texttt{cor\_03.jpg} du dossier illustrent le résultat attendu.
\end{corrigebox}
